\section{Смежные работы}
\label{sec:related}

Разнесение prefill и генерации по разным машинам хорошо изучено внутри дата-центра. DistServe и Splitwise отдают
фазы разным GPU, чтобы они не мешали друг другу \cite{zhong2024distserve,patel2023splitwise}; Mooncake строит
обслуживание вокруг распределённого хранилища KV-кэша \cite{qin2024mooncake}. Все они держат обе фазы на сервере. PrfaaS выносит длинный prefill на отдельные кластеры и передаёт кэш между
дата-центрами по обычному Ethernet, опираясь, как и мы, на модели с гибридным вниманием и малым кэшем
\cite{qin2026prfaas}; генерация у него остаётся в дата-центре. PPD для каждого следующего хода диалога решает, читать ли добавочный prefill на узле prefill или на
узле генерации, где уже лежит кэш \cite{li2026ppd}; выбор места для добавочного prefill предложен там, и обе
стороны у PPD работают внутри сервера. SparKV готовит на устройстве переиспользуемый контекст: облачный KV-кэш
делится на блоки по токенам, слоям и головам, и планировщик решает, какие блоки скачать, а какие посчитать на
месте, совмещает передачу с вычислениями и подстраивается под сеть \cite{liu2026sparkv}. Многоходовые диалоги
SparKV называет среди применений, но согласования двух растущих копий между ходами и рекуррентного состояния не
описывает; его планирование передачи может дополнить наш протокол.
Предложенное изменение ds4 распределяет prefill по слоям между несколькими машинами и затем ведёт генерацию на одной
из них \cite{ds4pr401}. Мы выносим генерацию и удержание кэша на машину пользователя, связываем её с поставщиком
prefill через интернет и держим у сторон две копии растущей агентской сессии. Новое в нашей схеме
относится к жизни сессии между ходами: клиент шлёт поставщику идентификаторы токенов, поставщик догоняет свою копию,
журнал и контрольные точки держат копии согласованными, рекуррентное состояние передаётся целиком, а клиент
продолжает работу без сервера.

Перенос состояния между моделями изучают как способ ускорить многоагентные системы. NVIDIA переносит кэш между
моделями одного семейства линейной картой \cite{heo2026crossmodel}; Cache-to-Cache проецирует кэш одной модели в
пространство другой нейросетью \cite{fu2025c2c}; DroidSpeak переиспользует кэш между дообученными вариантами одной
архитектуры и пересчитывает слои, где качество страдает \cite{liu2024droidspeak}. Ближе всего к нашим разнородным
квантам QSpec: две схемы квантизации одной модели делят один кэш при спекулятивном декодировании, и точная схема
переписывает кэш принятых токенов своими значениями \cite{zhao2024qspec}. Качество долгоживущей смешанной памяти,
как в нашей схеме, QSpec поэтому не подтверждает (ждёт нашей проверки). В нашей схеме кэш делят разные машины, и одна из сторон может не
доверять другой.

Проверяемый инференс через хэши промежуточных активаций предложен в TOPLOC \cite{ong2025toploc}. Атаки на KV-кэш,
которые меняют генерацию без изменения промпта, описаны в \cite{ganesh2025historyswap,zhu2026r2ecap,zhang2026hijackkv};
раздел~\ref{sec:security} показывает, какие проверки против них даёт гибридная схема.
